\section{Interface Boundaries and Recommendations}
\label{sec:root-causes}

Table~\ref{tab:root-causes} presents representative, not exhaustive, evidence chains linking a corpus clause or requirement identifier to a source/output observation and artifact-level interface mismatch. It does not assign blame or map every behavior to one permissive clause; defects, priorities, package-manager evidence, and configuration remain competing or interacting explanations.

To make these recommendations implementable, we propose a cumulative three-level defensive conformance-check suite. This remains a design proposal derived from our findings; implementation and evaluation are future work.

\noindent\textbf{Level 1: Identity and provenance.} Profile-specific schema and field-presence rules can check machine-readable supplier/author, component name and version, identifier, timestamp, and tool identity. Anchors include CycloneDX \field{metadata.supplier}, \field{metadata.timestamp}, \field{metadata.tools}, and component identity fields, and SPDX \field{CreationInfo.createdBy}, \field{createdUsing}, \field{packageVersion}, and \field{suppliedBy}.

\noindent\textbf{Level 2: Evidence and coverage.} Profile rules can combine producer declarations with checks for evidence sources, filtering policy, dependency relationships, coverage status, and absence reasons. CycloneDX \field{dependencies}, \field{scope}, and \field{compositions}, together with SPDX relationships and declared profiles, provide machine-readable anchors; evidence-source, filtering, and absence declarations require profile rules or extensions where no standard field exists.

\noindent\textbf{Level 3: Defensive-task sufficiency.} Task-specific checks can evaluate vulnerability/VEX context and build/formulation evidence through CycloneDX \field{vulnerabilities[].analysis} and \field{formulation} or SPDX Security and Build profiles. They then compare the artifact's claims with external manifests, lockfiles, builds, containers, or runtime evidence for vulnerability response, compliance, or incident triage.

Schemas can check the machine-expressible structure at each level, but they cannot establish the truth of producer declarations. The levels therefore keep specification completeness, generator-evidence completeness, and security-use-case sufficiency separate. Passing a lower level does not establish a higher one, and graph depth remains a diagnostic rather than a stand-alone completeness measure.
